\subsection{Adjoints and Orthogonal Operators}

% **Role:** linear algebra with metric and angle structure.
%
% Adjoints, conjugate transpose, self-adjoint maps, unitary/orthogonal maps, isometries, Householder reflections.

Adjoints express dual action in inner-product language. They become conjugate transposes in orthonormal coordinates and lead to self-adjoint maps, projections, isometries, and unitary or orthogonal operators.

\begin{definition} \label{def:adjoint-linear-map} \hlt{Adjoint of a Linear Map}\\
Let $\mathcal{V}$ and $\mathcal{W}$ be inner-product spaces and let $T\in\LL(\mathcal{V},\mathcal{W})$.\\
A function $S:\mathcal{W}\to\mathcal{V}$ is an \hlt{adjoint} of $T$ if $\innerprod[\mathcal{W}]{T\mathbf{v},\mathbf{w}} =\innerprod[\mathcal{V}]{\mathbf{v},S\mathbf{w}}$, $\mathbf{v}\in\mathcal{V},\ \mathbf{w}\in\mathcal{W}$.
\end{definition}

\begin{lemma} \label{lem:adjoint-unique-if-exists} \hlt{Uniqueness of Adjoints}\\
A linear map has at most one adjoint.
\begin{proof}
Suppose $S_1,S_2:\mathcal{W}\to\mathcal{V}$ are adjoints of $T$. For each $\mathbf{w}\in\mathcal{W}$ and every $\mathbf{v}\in\mathcal{V}$, $\innerprod[\mathcal{V}]{\mathbf{v},S_1\mathbf{w}-S_2\mathbf{w}}=0$.\\
Taking $\mathbf{v}=S_1\mathbf{w}-S_2\mathbf{w}$ gives $S_1\mathbf{w}=S_2\mathbf{w}$. Hence $S_1=S_2$.\\
\hlt{Structural Commentary}: Same tests force candidate adjoints to agree.\\
\hlt{Key Insight}: Inner products separate vectors.
\end{proof}
\end{lemma}

\begin{theorem} \label{thm:finite-dimensional-adjoint-exists} \hlt{Existence of the Adjoint}\\
Let $\mathcal{V}$ be a finite-dimensional inner-product space, let $\mathcal{W}$ be an inner-product space, and let $T\in\LL(\mathcal{V},\mathcal{W})$. Then $T$ has a unique adjoint $T^*:\mathcal{W}\to\mathcal{V}$, and $T^*$ is linear.
\begin{proof}
For each $\mathbf{w}\in\mathcal{W}$, the map $\varphi_{\mathbf{w}}:\mathcal{V}\to\mathbb{F}$ defined by $\varphi_{\mathbf{w}}(\mathbf{v})=\innerprod[\mathcal{W}]{T\mathbf{v},\mathbf{w}}$ is a linear functional. By Theorem \ref{thm:riesz-representation-finite-dimensional}, there is a unique vector $T^*\mathbf{w}\in\mathcal{V}$ such that $\innerprod[\mathcal{W}]{T\mathbf{v},\mathbf{w}} =\innerprod[\mathcal{V}]{\mathbf{v},T^*\mathbf{w}}$ for every $\mathbf{v}\in\mathcal{V}$.\\
Thus $T^*$ is an adjoint, and it is unique by Lemma \ref{lem:adjoint-unique-if-exists}. For $\mathbf{w}_1,\mathbf{w}_2\in\mathcal{W}$ and $a\in\mathbb{F}$,
\begin{align}
\innerprod[\mathcal{V}]{\mathbf{v},T^*(a\mathbf{w}_1+\mathbf{w}_2)}
=\innerprod[\mathcal{W}]{T\mathbf{v},a\mathbf{w}_1+\mathbf{w}_2}
=\overline{a}\innerprod[\mathcal{W}]{T\mathbf{v},\mathbf{w}_1}+\innerprod[\mathcal{W}]{T\mathbf{v},\mathbf{w}_2}
=\innerprod[\mathcal{V}]{\mathbf{v},aT^*\mathbf{w}_1+T^*\mathbf{w}_2}. \nonumber
\end{align}
for every $\mathbf{v}\in\mathcal{V}$. Taking the difference of the two second-slot vectors and then setting $\mathbf{v}$ equal to that difference gives $T^*(a\mathbf{w}_1+\mathbf{w}_2)=aT^*\mathbf{w}_1+T^*\mathbf{w}_2$.\\
\hlt{Structural Commentary}: Each $\mathbf{w}$ gives a Riesz-represented functional on $\mathcal{V}$.\\
\hlt{Key Insight}: The adjoint is pointwise Riesz representation for $T$.
\end{proof}
\end{theorem}

\begin{figure}[H]
  \centering
  \begin{tikzpicture}[
  >=Stealth,
  every node/.style={font=\small},
  space/.style={draw=gray!45, rounded corners=3pt, minimum width=3.6cm, minimum height=2.4cm},
  vec/.style={->, blue!75!black, line width=0.85pt},
  adj/.style={->, teal!75!black, line width=0.85pt},
  map/.style={->, gray!70, line width=0.75pt}
]
  \node[space] (V) at (0,0) {};
  \node[anchor=north west] at (-1.6,1.05) {$\mathcal{V}$};
  \node[space] (W) at (5.2,0) {};
  \node[anchor=north west] at (3.6,1.05) {$\mathcal{W}$};

  \coordinate (Ov) at (-0.95,-0.65);
  \coordinate (v) at (0.55,0.55);
  \coordinate (Tsw) at (0.85,-0.25);
  \draw[vec] (Ov) -- (v) node[midway, above left] {$\mathbf{v}$};
  \draw[adj] (Ov) -- (Tsw) node[midway, below] {$T^*\mathbf{w}$};

  \coordinate (Ow) at (4.25,-0.65);
  \coordinate (Tv) at (5.85,0.55);
  \coordinate (w) at (5.95,-0.25);
  \draw[vec] (Ow) -- (Tv) node[midway, above left] {$T\mathbf{v}$};
  \draw[adj] (Ow) -- (w) node[midway, below] {$\mathbf{w}$};

  \draw[map] (1.9,0.45) -- (3.25,0.45) node[midway, above] {$T$};
  \draw[map] (3.25,-0.45) -- (1.9,-0.45) node[midway, below] {$T^*$};

  \node[draw=gray!45, rounded corners=3pt, align=center, inner sep=5pt] at (2.6,-1.85)
    {$\innerprod[\mathcal{W}]{T\mathbf{v},\mathbf{w}}
      =\innerprod[\mathcal{V}]{\mathbf{v},T^*\mathbf{w}}$};
\end{tikzpicture}

  \caption{Adjoint relation: $\innerprod[\mathcal{W}]{T\mathbf{v},\mathbf{w}}=\innerprod[\mathcal{V}]{\mathbf{v},T^*\mathbf{w}}$.}
  \label{fig:adjoint-inner-product-bridge}
\end{figure}

\begin{definition} \label{def:conjugate-transpose} \hlt{Conjugate Transpose}\\
For $A=[a_{ij}]\in\operatorname{Mat}(m,n,\mathbb{F})$, the \hlt{conjugate transpose} of $A$ is $A^*=[\overline{a_{ji}}]\in\operatorname{Mat}(n,m,\mathbb{F})$.\\
Over $\mathbb{R}$, this is the ordinary transpose.
\end{definition}

\begin{theorem} \label{thm:matrix-of-adjoint-is-conjugate-transpose} \hlt{Matrix of the Adjoint}\\
Let $\mathcal{V}$ and $\mathcal{W}$ be finite-dimensional inner-product spaces. Let $\beta=(\mathbf{e}_1,\ldots,\mathbf{e}_n)$ be an orthonormal basis of $\mathcal{V}$ and let $\gamma=(\mathbf{f}_1,\ldots,\mathbf{f}_m)$ be an orthonormal basis of $\mathcal{W}$. If $T\in\LL(\mathcal{V},\mathcal{W})$, then $[T^*]_{\beta\leftarrow\gamma} =\bigl([T]_{\gamma\leftarrow\beta}\bigr)^*$.
\begin{proof}
Let $A=[T]_{\gamma\leftarrow\beta}$ and $B=[T^*]_{\beta\leftarrow\gamma}$.\\
By Theorem \ref{thm:matrix-entries-from-orthonormal-codomain}, $B_{i,j}=\innerprod[\mathcal{V}]{T^*\mathbf{f}_j,\mathbf{e}_i} =\overline{\innerprod[\mathcal{V}]{\mathbf{e}_i,T^*\mathbf{f}_j}} =\overline{\innerprod[\mathcal{W}]{T\mathbf{e}_i,\mathbf{f}_j}} =\overline{A_{j,i}}$. Thus $B=A^*$.\\
\hlt{Structural Commentary}: Entries swap by moving $T$ across the inner product.\\
\hlt{Key Insight}: In orthonormal coordinates, taking adjoints is exactly taking conjugate transposes.
\end{proof}
\end{theorem}

\begin{proposition} \label{prop:basic-adjoint-identities} \hlt{Basic Adjoint Identities}\\
For finite-dimensional inner-product spaces and composable linear maps:
\begin{enumerate}[label=\roman*.]
\setlength{\itemsep}{0pt}
\item $(aS+T)^*=\overline{a}S^*+T^*$ for $S,T\in\LL(\mathcal{V},\mathcal{W})$ and $a\in\mathbb{F}$;
\item $(SR)^*=R^*S^*$ for $R\in\LL(\mathcal{U},\mathcal{V})$ and $S\in\LL(\mathcal{V},\mathcal{W})$;
\item $(T^*)^*=T$;
\item $I_{\mathcal{V}}^*=I_{\mathcal{V}}$.
\end{enumerate}
\begin{proof}
Verify the defining identity; uniqueness gives equality.
\begin{enumerate}[label=(\roman*)]
\setlength{\itemsep}{0pt}
\item For $\mathbf{v}\in\mathcal{V}$, $\mathbf{w}\in\mathcal{W}$, $\innerprod{(aS+T)\mathbf{v},\mathbf{w}} =a\innerprod{S\mathbf{v},\mathbf{w}}+\innerprod{T\mathbf{v},\mathbf{w}} =\innerprod{\mathbf{v},\overline{a}S^*\mathbf{w}+T^*\mathbf{w}}$, so $(aS+T)^*=\overline{a}S^*+T^*$.
\item For $\mathbf{u}\in\mathcal{U}$ and $\mathbf{w}\in\mathcal{W}$, $\innerprod{SR\mathbf{u},\mathbf{w}} =\innerprod{R\mathbf{u},S^*\mathbf{w}} =\innerprod{\mathbf{u},R^*S^*\mathbf{w}}$, so $(SR)^*=R^*S^*$.
\item Conjugating $\innerprod{T\mathbf{v},\mathbf{w}}=\innerprod{\mathbf{v},T^*\mathbf{w}}$ gives $\innerprod{T^*\mathbf{w},\mathbf{v}}=\innerprod{\mathbf{w},T\mathbf{v}}$, so $T$ is the adjoint of $T^*$.
\item Since $\innerprod{I\mathbf{v},\mathbf{w}}=\innerprod{\mathbf{v},I\mathbf{w}}$, the identity map is its own adjoint.
\end{enumerate}
\hlt{Structural Commentary}: Put both sides in one inner-product slot, then use uniqueness.\\
\hlt{Key Insight}: Like matrix transposes, adjoints reverse products and conjugate scalars.
\end{proof}
\end{proposition}

\begin{proposition} \label{prop:adjoint-kernel-range-orthogonality} \hlt{Adjoint Kernel-Range Orthogonality}\\
Let $\mathcal{V}$ and $\mathcal{W}$ be finite-dimensional inner-product spaces and let $T\in\LL(\mathcal{V},\mathcal{W})$.\\
Then $\ker(T^*)=(\operatorname{range}(T))^{\perp}$, $\operatorname{range}(T^*)=(\ker(T))^{\perp}$, and $\ker(T)=(\operatorname{range}(T^*))^{\perp}$, $\operatorname{range}(T)=(\ker(T^*))^{\perp}$.\\
Equivalently, the nullspace of an adjoint is the orthogonal complement of the range of the original map.
\begin{proof}
For $\mathbf{w}\in\mathcal{W}$,
\begin{align}
\mathbf{w}\in\ker(T^*)
&\Longleftrightarrow \innerprod[\mathcal{V}]{\mathbf{v},T^*\mathbf{w}}=0\quad\text{for every }\mathbf{v}\in\mathcal{V} \nonumber\\
&\Longleftrightarrow \innerprod[\mathcal{W}]{T\mathbf{v},\mathbf{w}}=0\quad\text{for every }\mathbf{v}\in\mathcal{V} \nonumber\\
&\Longleftrightarrow \mathbf{w}\in(\operatorname{range}(T))^{\perp}.
\nonumber
\end{align}
Applying this identity to $T^*$ gives $\ker(T)=(\operatorname{range}(T^*))^{\perp}$.\\
Since all spaces are finite-dimensional, the ranges and kernels are finite-dimensional subspaces.\\
Taking orthogonal complements and using Corollary \ref{cor:double-orthogonal-complement} gives
\begin{align}
(\ker(T^*))^{\perp}=((\operatorname{range}(T))^{\perp})^{\perp}=\operatorname{range}(T),
\qquad
(\ker(T))^{\perp}=((\operatorname{range}(T^*))^{\perp})^{\perp}=\operatorname{range}(T^*) \nonumber
\end{align}
\hlt{Structural Commentary}: Adjoints trade ranges for kernels, and orthogonal complements reverse that trade.\\
\hlt{Key Insight}: A vector is killed by $T^*$ exactly when it is orthogonal to everything $T$ can produce.
\end{proof}
\end{proposition}

\begin{corollary} \label{cor:matrix-fundamental-subspaces-orthogonal-complements} \hlt{Matrix Fundamental Subspaces as Orthogonal Complements}\\
Let $A\in\operatorname{Mat}(m,n,\mathbb{F})$ act as $A:\mathbb{F}^n\to\mathbb{F}^m$ with standard inner products.\\
Then $\ker(A)=\operatorname{col}(A^*)^{\perp}$ and $\operatorname{col}(A^*)&=\ker(A)^{\perp}$; $\ker(A^*)=\operatorname{col}(A)^{\perp}$ and $\operatorname{col}(A)&=\ker(A^*)^{\perp}$.\\
Over $\mathbb{R}$, this says $\ker(A)\perp\operatorname{row}(A)$ in $\mathbb{R}^n$ and $\ker(A^t)\perp\operatorname{col}(A)$ in $\mathbb{R}^m$, with both pairs complementary.
\begin{proof}
By Theorem \ref{thm:matrix-of-adjoint-is-conjugate-transpose}, the adjoint of $A$ is represented by $A^*$.\\
Proposition \ref{prop:adjoint-kernel-range-orthogonality} gives $\ker(A^*)=\operatorname{col}(A)^{\perp}$ and $\ker(A)=\operatorname{col}(A^*)^{\perp}$. Taking orthogonal complements and using Corollary \ref{cor:double-orthogonal-complement} gives the two column-space identities. Over $\mathbb{R}$, $A^*=A^t$ and $\operatorname{col}(A^t)=\operatorname{row}(A)$.\\
\hlt{Structural Commentary}: The four fundamental subspaces form two orthogonal-complement pairs.\\
\hlt{Key Insight}: Adjoints turn missing directions into kernels.
\end{proof}
\end{corollary}

\begin{definition} \label{def:self-adjoint-hermitian} \hlt{Self-Adjoint and Hermitian}\\
An operator $T\in\LL(\mathcal{V})$ on a finite-dimensional inner-product space is \hlt{self-adjoint} if $T^*=T$. A square matrix $A$ is \hlt{Hermitian} if $A^*=A$. In orthonormal coordinates, these are the same condition by Theorem \ref{thm:matrix-of-adjoint-is-conjugate-transpose}.
\end{definition}

\begin{lemma} \label{lem:complex-zero-quadratic-form-implies-zero-operator} \hlt{Zero Complex Quadratic Form}\\
Let $\mathcal{V}$ be a finite-dimensional complex inner-product space and let $T\in\LL(\mathcal{V})$.\\
If $\innerprod{T\mathbf{v},\mathbf{v}}=0$ for every $\mathbf{v}\in\mathcal{V}$, then $T=0$.
\begin{proof}
Set $Q(\mathbf{x})=\innerprod{T\mathbf{x},\mathbf{x}}$.\\
For $\mathbf{u},\mathbf{w}\in\mathcal{V}$, expansion gives the polarization identity for the sesquilinear form $(\mathbf{u},\mathbf{w})\mapsto\innerprod{T\mathbf{u},\mathbf{w}}$:
\begin{align}
\innerprod{T\mathbf{u},\mathbf{w}} =\frac{Q(\mathbf{u}+\mathbf{w})-Q(\mathbf{u}-\mathbf{w})}{4} +i\,\frac{Q(\mathbf{u}+i\mathbf{w})-Q(\mathbf{u}-i\mathbf{w})}{4}. \nonumber
\end{align}
The hypothesis makes every term on the right zero, so $\innerprod{T\mathbf{u},\mathbf{w}}=0$ for all $\mathbf{u},\mathbf{w}$.\\
Taking $\mathbf{w}=T\mathbf{u}$ gives $\norm{T\mathbf{u}}^2=0$, hence $T\mathbf{u}=\mathbf{0}$ for every $\mathbf{u}$.\\
\hlt{Structural Commentary}: Over $\mathbb{C}$, diagonal values of a sesquilinear form determine all mixed values.\\
\hlt{Key Insight}: Polarization turns the one-vector condition into a two-vector condition.
\end{proof}
\end{lemma}

\begin{proposition} \label{prop:complex-quadratic-form-self-adjoint-characterization} \hlt{Self-Adjointness from Real Quadratic Values}\\
Let $\mathcal{V}$ be a finite-dimensional complex inner-product space and let $T\in\LL(\mathcal{V})$.\\
Then $T$ is self-adjoint if and only if $\innerprod{T\mathbf{v},\mathbf{v}}\in\mathbb{R}$ for every $\mathbf{v}\in\mathcal{V}$.
\begin{proof}
For each $\mathbf{v}$, $\innerprod{(T-T^*)\mathbf{v},\mathbf{v}} =\innerprod{T\mathbf{v},\mathbf{v}}-\innerprod{T^*\mathbf{v},\mathbf{v}} =\innerprod{T\mathbf{v},\mathbf{v}}-\overline{\innerprod{T\mathbf{v},\mathbf{v}}}$.\\
If all values $\innerprod{T\mathbf{v},\mathbf{v}}$ are real, this expression is zero for every $\mathbf{v}$, and Lemma \ref{lem:complex-zero-quadratic-form-implies-zero-operator} gives $T-T^*=0$.\\
Conversely, if $T=T^*$, then $\innerprod{T\mathbf{v},\mathbf{v}} =\innerprod{\mathbf{v},T\mathbf{v}} =\overline{\innerprod{T\mathbf{v},\mathbf{v}}}$, so $\innerprod{T\mathbf{v},\mathbf{v}}$ is real.\\
\hlt{Structural Commentary}: The imaginary part of $\innerprod{T\mathbf{v},\mathbf{v}}$ is exactly the diagonal trace of $T-T^*$.\\
\hlt{Key Insight}: Over $\mathbb{C}$, real quadratic values force the operator to equal its adjoint.
\end{proof}
\end{proposition}

\begin{proposition} \label{prop:self-adjoint-zero-quadratic-form-implies-zero} \hlt{Zero Quadratic Form for Self-Adjoint Operators}\\
Let $T$ be a self-adjoint operator on a finite-dimensional inner-product space $\mathcal{V}$.\\
If $\innerprod{T\mathbf{v},\mathbf{v}}=0$ for every $\mathbf{v}\in\mathcal{V}$, then $T=0$.
\begin{proof}
If $\mathbb{F}=\mathbb{C}$, this is Lemma \ref{lem:complex-zero-quadratic-form-implies-zero-operator}. Assume $\mathbb{F}=\mathbb{R}$.\\
For $\mathbf{u},\mathbf{w}\in\mathcal{V}$, self-adjointness gives $\innerprod{T\mathbf{w},\mathbf{u}}=\innerprod{\mathbf{w},T\mathbf{u}}=\innerprod{T\mathbf{u},\mathbf{w}}$, so the real polarization identity gives
\begin{align}
\innerprod{T\mathbf{u},\mathbf{w}}
=\frac{\innerprod{T(\mathbf{u}+\mathbf{w}),\mathbf{u}+\mathbf{w}}-\innerprod{T(\mathbf{u}-\mathbf{w}),\mathbf{u}-\mathbf{w}}}{4}=0.
\nonumber
\end{align}
Taking $\mathbf{w}=T\mathbf{u}$ gives $T\mathbf{u}=\mathbf{0}$ for every $\mathbf{u}$.\\
\hlt{Structural Commentary}: Self-adjointness supplies the missing symmetry needed for real polarization.\\
\hlt{Key Insight}: A self-adjoint operator is determined by its quadratic form.
\end{proof}
\end{proposition}

\begin{definition} \label{def:orthogonal-projection-matrix} \hlt{Orthogonal Projection Matrix}\\
A square matrix $P\in\operatorname{Mat}(n,n,\mathbb{F})$ is an \hlt{orthogonal projection matrix} if $P^2=P$ and $P^*=P$.
\end{definition}

\begin{theorem} \label{thm:orthogonal-projection-self-adjoint-idempotent} \hlt{Orthogonal Projections as Self-Adjoint Idempotents}\\
Let $\mathcal{V}$ be finite-dimensional and let $P\in\LL(\mathcal{V})$.\\
Then $P$ is the orthogonal projection onto $\operatorname{range}(P)$ if and only if $P^2=P$ and $P^*=P$.
\begin{proof}
\begin{enumerate}[label=\roman*.]
\setlength{\itemsep}{0pt}
\item[$(\Rightarrow)$] If $P=P_{\mathcal{U}}$ for $\mathcal{U}=\operatorname{range}(P)$, then Proposition \ref{prop:orthogonal-projection-properties} gives $P^2=P$. For $\mathbf{v}=\mathbf{u}+\mathbf{w}$ and $\mathbf{x}=\mathbf{a}+\mathbf{b}$ with $\mathbf{u},\mathbf{a}\in\mathcal{U}$ and $\mathbf{w},\mathbf{b}\in\mathcal{U}^{\perp}$, $\innerprod{P\mathbf{v},\mathbf{x}}=\innerprod{\mathbf{u},\mathbf{a}+\mathbf{b}}=\innerprod{\mathbf{u},\mathbf{a}} =\innerprod{\mathbf{u}+\mathbf{w},\mathbf{a}}=\innerprod{\mathbf{v},P\mathbf{x}}$, so $P^*=P$.
\item[$(\Leftarrow)$] Suppose $P^2=P$ and $P^*=P$. For each $\mathbf{v}$, $\mathbf{v}=P\mathbf{v}+(\mathbf{v}-P\mathbf{v})$ with $P\mathbf{v}\in\operatorname{range}(P)$ and $P(\mathbf{v}-P\mathbf{v})=\mathbf{0}$.\\
Thus the residual lies in $\ker(P)$. Since $P=P^*$, Proposition \ref{prop:adjoint-kernel-range-orthogonality} gives $\ker(P)=(\operatorname{range}(P))^{\perp}$.\\
Hence every $\mathbf{v}$ has the orthogonal decomposition $\mathbf{v}=P\mathbf{v}+(\mathbf{v}-P\mathbf{v})$ with first part in $\operatorname{range}(P)$, so $P$ is exactly the orthogonal projection onto $\operatorname{range}(P)$.
\end{enumerate}
\hlt{Structural Commentary}: Idempotence splits by range and kernel; self-adjointness makes the split orthogonal.\\
\hlt{Key Insight}: An idempotent is orthogonal exactly when $\ker(P)=(\operatorname{range}(P))^{\perp}$.
\end{proof}
\end{theorem}

\begin{theorem} \label{thm:orthogonal-projection-matrix-unitary-coordinate-form} \hlt{Unitary Coordinate Form of an Orthogonal Projection}\\
Let $P\in\operatorname{Mat}(n,n,\mathbb{F})$ be an orthogonal projection matrix and let $r=\operatorname{rank}(P)$. Then $P$ is unitarily similar to $I_r\oplus0_{n-r}$. Equivalently, there is a unitary matrix $U$ such that
\begin{align}
U^*PU=
\begin{bmatrix}
I_r & 0\\
0 & 0_{n-r}
\end{bmatrix}. \nonumber
\end{align}
\begin{proof}
Let $\mathcal{U}=\operatorname{col}(P)$. Identifying $P$ with its induced operator on $\mathbb{F}^n$, Theorem \ref{thm:orthogonal-projection-self-adjoint-idempotent} says $P$ is the orthogonal projection onto $\mathcal{U}$, so $P$ is the identity on $\mathcal{U}$ and zero on $\mathcal{U}^{\perp}$. Choose an orthonormal basis $\mathbf{u}_1,\ldots,\mathbf{u}_r$ of $\mathcal{U}$ and extend it to an orthonormal basis $\mathbf{u}_1,\ldots,\mathbf{u}_n$ of $\mathbb{F}^n$. Let $U=[\mathbf{u}_1\ \cdots\ \mathbf{u}_n]$.\\
Then $U$ is unitary. For $j\leq r$, $P\mathbf{u}_j=\mathbf{u}_j$, and for $j>r$, $\mathbf{u}_j\in\mathcal{U}^{\perp}$ so $P\mathbf{u}_j=\mathbf{0}$. Therefore the $j$th column of $U^*PU$ is the coordinate column of $P\mathbf{u}_j$ in this orthonormal basis, giving the displayed block diagonal form.\\
\hlt{Structural Commentary}: Orthogonal projection splits the space into range and orthogonal kernel.\\
\hlt{Key Insight}: In a basis adapted to $\operatorname{col}(P)\oplus\operatorname{col}(P)^{\perp}$, projection keeps the first coordinates and kills the rest.
\end{proof}
\end{theorem}

\begin{figure}[H]
  \centering
  \begin{tikzpicture}[
  x=0.88cm,
  y=0.88cm,
  >=Stealth,
  every node/.style={font=\small},
  sub/.style={draw=blue!60!black, fill=blue!7, rounded corners=3pt, minimum width=2.07cm, minimum height=0.64cm, align=center},
  mat/.style={draw=gray!60, fill=gray!5, rounded corners=2pt, minimum width=1.90cm, minimum height=1.58cm},
  arrow/.style={->, line width=0.85pt, blue!70!black}
]
  \node[sub] (C) at (0,0.72) {$\operatorname{col}(P)$};
  \node[sub, fill=red!6, draw=red!55!black] (N) at (0,-0.72) {$\operatorname{col}(P)^\perp$};

  \node[mat] (B) at (4.1,0) {};
  \draw[gray!55] ($(B.center)+(-0.86,0)$) -- ($(B.center)+(0.86,0)$);
  \draw[gray!55] ($(B.center)+(0,-0.68)$) -- ($(B.center)+(0,0.68)$);
  \node[blue!70!black] at ($(B.center)+(-0.43,0.43)$) {$I_r$};
  \node[red!60!black] at ($(B.center)+(0.43,-0.43)$) {$0_{n-r}$};
  \node[gray!65] at ($(B.center)+(0.43,0.43)$) {$0$};
  \node[gray!65] at ($(B.center)+(-0.43,-0.43)$) {$0$};

  \draw[arrow] (C.east) -- ($(B.west)+(0,0.43)$);
  \draw[arrow, red!60!black] (N.east) -- ($(B.west)+(0,-0.43)$);
\end{tikzpicture}

  \caption{Adapted basis: $U^*PU=I_r\oplus0_{n-r}$, identity on $\operatorname{col}(P)$ and zero on $\operatorname{col}(P)^\perp$.}
  \label{fig:orthogonal-projection-unitary-block}
\end{figure}

\begin{definition} \label{def:isometry-linear-map} \hlt{Isometry}\\
Let $\mathcal{V}$ be an inner-product space. An operator $T\in\LL(\mathcal{V})$ is an \hlt{isometry} if $\norm{T\mathbf{v}}=\norm{\mathbf{v}}$ for every $\mathbf{v}\in\mathcal{V}$.
\end{definition}

\begin{proposition} \label{prop:basic-isometry-properties} \hlt{Basic Properties of Isometries}\\
Let $\mathcal{V}$ be an inner-product space and let $S,T\in\LL(\mathcal{V})$ be isometries.
\begin{enumerate}[label=\roman*.]
\setlength{\itemsep}{0pt}
\item $ST$ is an isometry;
\item $T$ is injective;
\item if $\mathcal{V}$ is finite-dimensional, then $T$ is invertible and $T^{-1}$ is an isometry.
\end{enumerate}
\begin{proof}
\begin{enumerate}[label=(\roman*)]
\setlength{\itemsep}{0pt}
\item For every $\mathbf{v}$, $\norm{ST\mathbf{v}}=\norm{T\mathbf{v}}=\norm{\mathbf{v}}$.
\item If $T\mathbf{v}=\mathbf{0}$, then $\norm{\mathbf{v}}=\norm{T\mathbf{v}}=0$, so $\mathbf{v}=\mathbf{0}$. Thus $\ker(T)=\{\mathbf{0}\}$.
\item In finite dimensions, injectivity implies invertibility by Theorem \ref{thm:finite-dimensional-operator-invertible-equivalences}.\\
If $\mathbf{u}=T\mathbf{v}$, then $\norm{T^{-1}\mathbf{u}}=\norm{\mathbf{v}}=\norm{T\mathbf{v}}=\norm{\mathbf{u}}$.
\end{enumerate}
\hlt{Structural Commentary}: Isometries preserve length, so they cannot collapse nonzero vectors.\\
\hlt{Key Insight}: In finite dimensions, length preservation already forces an invertible geometry-preserving map.
\end{proof}
\end{proposition}

\begin{theorem} \label{thm:isometry-iff-inner-product-preserving} \hlt{Isometries Preserve Inner Products}\\
Let $\mathcal{V}$ be an inner-product space and let $T\in\LL(\mathcal{V})$.\\
Then $T$ is an isometry if and only if $\innerprod{T\mathbf{u},T\mathbf{v}}=\innerprod{\mathbf{u},\mathbf{v}}$ for every $\mathbf{u},\mathbf{v}\in\mathcal{V}$.
\begin{proof}
\begin{enumerate}[label=\roman*.]
\setlength{\itemsep}{0pt}
\item[$(\Rightarrow)$] Suppose $T$ is an isometry. Then every norm appearing in the polarization identities of Theorem \ref{thm:polarization-identities} is unchanged after applying $T$. Therefore $\innerprod{T\mathbf{u},T\mathbf{v}}=\innerprod{\mathbf{u},\mathbf{v}}$.
\item[$(\Leftarrow)$] Suppose $T$ preserves inner products. Then $\norm{T\mathbf{u}}^2=\innerprod{T\mathbf{u},T\mathbf{u}}=\innerprod{\mathbf{u},\mathbf{u}}=\norm{\mathbf{u}}^2$, so $T$ is an isometry.
\end{enumerate}
\hlt{Structural Commentary}: Polarization converts preservation of norms into preservation of inner products.\\
\hlt{Key Insight}: Length preservation for all vectors also preserves angle information.
\end{proof}
\end{theorem}

\begin{theorem} \label{thm:finite-dimensional-isometry-iff-adjoint-product-identity} \hlt{Isometries and the Adjoint Product}\\
Let $\mathcal{V}$ be finite-dimensional and let $T\in\LL(\mathcal{V})$. Then $T$ is an isometry if and only if $T^*T=I_{\mathcal{V}}$.
\begin{proof}
\begin{enumerate}[label=\roman*.]
\setlength{\itemsep}{0pt}
\item[$(\Rightarrow)$] Suppose $T$ is an isometry. By Theorem \ref{thm:isometry-iff-inner-product-preserving}, $\innerprod{T^*T\mathbf{u},\mathbf{v}}=\innerprod{T\mathbf{u},T\mathbf{v}}=\innerprod{\mathbf{u},\mathbf{v}}$ for all $\mathbf{u},\mathbf{v}$.\\
Thus $\innerprod{(T^*T-I_{\mathcal{V}})\mathbf{u},\mathbf{v}}=0$ for every $\mathbf{v}$, so $T^*T\mathbf{u}=\mathbf{u}$ for every $\mathbf{u}$.
\item[$(\Leftarrow)$] Suppose $T^*T=I_{\mathcal{V}}$. Then $\norm{T\mathbf{u}}^2 =\innerprod{T\mathbf{u},T\mathbf{u}} =\innerprod{T^*T\mathbf{u},\mathbf{u}} =\norm{\mathbf{u}}^2$ for every $\mathbf{u}$, so $T$ is an isometry.
\end{enumerate}
\hlt{Structural Commentary}: The operator $T^*T$ records inner products after applying $T$.\\
\hlt{Key Insight}: $T$ is an isometry exactly when this Gram operator is the identity.
\end{proof}
\end{theorem}

\begin{definition} \label{def:unitary-operator-and-matrix} \hlt{Unitary and Orthogonal}\\
An operator $U\in\LL(\mathcal{V})$ on a finite-dimensional inner-product space is \hlt{unitary} if $U^*U=UU^*=I_{\mathcal{V}}$.\\
A matrix $U\in\operatorname{Mat}(n,n,\mathbb{F})$ is \hlt{unitary} if $U^*U=I_n$. If $\mathbb{F}=\mathbb{R}$, a unitary matrix is called \hlt{orthogonal}.
\end{definition}

\begin{figure}[H]
  \centering
  \begin{tikzpicture}[
  >=Stealth,
  every node/.style={font=\small},
  circleline/.style={gray!50, line width=0.65pt},
  veca/.style={->, blue!75!black, line width=0.95pt},
  vecb/.style={->, teal!70!black, line width=0.95pt},
  map/.style={->, gray!70, line width=0.8pt}
]
  \coordinate (O1) at (0,0);
  \coordinate (O2) at (5.0,0);
  \coordinate (Ua) at ($(O2)+(0.72,0.72)$);
  \coordinate (Ub) at ($(O2)+(-0.72,0.72)$);
  \draw[circleline] (O1) circle (1.15);
  \draw[circleline] (O2) circle (1.15);

  \draw[veca] (O1) -- (1.0,0) node[below] {$\mathbf{e}_1$};
  \draw[vecb] (O1) -- (0,1.0) node[left] {$\mathbf{e}_2$};
  \draw[veca] (O2) -- (Ua) node[above right] {$U\mathbf{e}_1$};
  \draw[vecb] (O2) -- (Ub) node[above left] {$U\mathbf{e}_2$};

  \draw[gray!55] (0.82,0) arc[start angle=0,end angle=90,radius=0.82];
  \draw[gray!55] ($(O2)+(0.58,0.58)$) arc[start angle=45,end angle=135,radius=0.82];
  \node at (0.55,0.55) {$\perp$};
  \node at ($(O2)+(0,0.78)$) {$\perp$};

  \draw[map] (1.65,0.15) -- (3.35,0.15) node[midway, above] {$U$};
\end{tikzpicture}

  \caption{Unitary invariance: $\norm{U\mathbf{x}}=\norm{\mathbf{x}}$ and $\innerprod{U\mathbf{x},U\mathbf{y}}=\innerprod{\mathbf{x},\mathbf{y}}$.}
  \label{fig:unitary-isometry-geometry}
\end{figure}

\begin{theorem} \label{thm:unitary-matrix-equivalent-conditions} \hlt{Equivalent Conditions for a Unitary Matrix}\\
Let $U=[\mathbf{u}_1\ \cdots\ \mathbf{u}_n]\in\operatorname{Mat}(n,n,\mathbb{F})$. The following are equivalent:
\begin{enumerate}[label=\roman*.]
\setlength{\itemsep}{0pt}
\item $U$ is unitary, i.e. $U^*U=I_n$;
\item the columns $\mathbf{u}_1,\ldots,\mathbf{u}_n$ are orthonormal in $\mathbb{F}^n$;
\item $UU^*=I_n$;
\item $U^*$ is unitary;
\item the rows of $U$ are orthonormal;
\item $U$ is invertible and $U^{-1}=U^*$;
\item $\norm{U\mathbf{x}}_2=\norm{\mathbf{x}}_2$ for every $\mathbf{x}\in\mathbb{F}^n$;
\item $\innerprod[\mathbb{F}^n]{U\mathbf{x},U\mathbf{y}}=\innerprod[\mathbb{F}^n]{\mathbf{x},\mathbf{y}}$ for every $\mathbf{x},\mathbf{y}\in\mathbb{F}^n$.
\end{enumerate}
\begin{proof}
\begin{enumerate}[label=(\roman*).]
\setlength{\itemsep}{0pt}
\item[$(i)$]$\Rightarrow (ii)$. Since $(U^*U)_{ij}=\mathbf{u}_i^*\mathbf{u}_j=\innerprod[\mathbb{F}^n]{\mathbf{u}_j,\mathbf{u}_i}$, equation $U^*U=I_n$ states that columns of $U$ are orthonormal.
\item[$(ii)$]$\Rightarrow (iii)$. Orthonormal columns give $U^*U=I_n$. A left inverse of a square matrix is also a right inverse by Theorem \ref{thm:square-one-sided-inverse-is-two-sided}, so $UU^*=I_n$.
\item[$(iii)$]$\Rightarrow (iv)$. The equation $UU^*=I_n$ is $(U^*)^*U^*=I_n$, so $U^*$ is unitary.
\item[$(iv)$]$\Rightarrow (v)$. Apply $(i)\Rightarrow(ii)$ to $U^*$; the columns of $U^*$ are the conjugates of the rows of $U$.
\item[$(v)$]$\Rightarrow (vi)$. Orthonormal rows give $UU^*=I_n$. Thus $U^*$ is the inverse of $U$, so $U^{-1}=U^*$.
\item[$(vi)$]$\Rightarrow (vii)$. If $U^{-1}=U^*$, then $U^*U=I_n$. Therefore $\norm{U\mathbf{x}}_2^2 =\innerprod[\mathbb{F}^n]{U\mathbf{x},U\mathbf{x}} =\innerprod[\mathbb{F}^n]{U^*U\mathbf{x},\mathbf{x}} =\norm{\mathbf{x}}_2^2$ for every $\mathbf{x}\in\mathbb{F}^n$. Since both norms are nonnegative, $\norm{U\mathbf{x}}_2=\norm{\mathbf{x}}_2$.
\item[$(vii)$]$\Rightarrow (viii)$. Apply Theorem \ref{thm:isometry-iff-inner-product-preserving} to $\mathbf{x}\mapsto U\mathbf{x}$.
\item[$(viii)$]$\Rightarrow (i)$. For every $\mathbf{x},\mathbf{y}\in\mathbb{F}^n$, $\innerprod[\mathbb{F}^n]{(U^*U-I_n)\mathbf{x},\mathbf{y}} =\innerprod[\mathbb{F}^n]{U\mathbf{x},U\mathbf{y}}-\innerprod[\mathbb{F}^n]{\mathbf{x},\mathbf{y}} =0$. Taking $\mathbf{y}=(U^*U-I_n)\mathbf{x}$ shows $(U^*U-I_n)\mathbf{x}=\mathbf{0}$ for every $\mathbf{x}$, so $U^*U=I_n$.
\end{enumerate}
\hlt{Structural Commentary}: A unitary matrix is an isometry written in orthonormal coordinates.\\
\hlt{Key Insight}: Columns, rows, inverses, norms, and inner products are different faces of the same condition.
\end{proof}
\end{theorem}

\begin{proposition} \label{prop:algebra-of-unitary-matrices} \hlt{Algebra of Unitary Matrices}\\
Let $U,V\in\operatorname{Mat}(n,n,\mathbb{F})$ and $W\in\operatorname{Mat}(m,m,\mathbb{F})$, where $V\oplus W$ denotes the block diagonal matrix with blocks $V$ and $W$.
\begin{enumerate}[label=\roman*.]
\setlength{\itemsep}{0pt}
\item If $U$ and $V$ are unitary, then $UV$ is unitary.
\item $V\oplus W$ is unitary if and only if $V$ and $W$ are unitary.
\end{enumerate}
\begin{proof}
\begin{enumerate}[label=(\roman*)]
\setlength{\itemsep}{0pt}
\item Compute $(UV)^*(UV)=V^*U^*UV=V^*V=I_n$.
\item Since $(V\oplus W)^*(V\oplus W)=V^*V\oplus W^*W$, this product equals $I_{n+m}$ exactly when $V^*V=I_n$ and $W^*W=I_m$.
\end{enumerate}
\hlt{Structural Commentary}: Unitary structure is stable under composition and blockwise orthogonal sums.\\
\hlt{Key Insight}: The equation $U^*U=I$ makes every verification a short adjoint calculation.
\end{proof}
\end{proposition}

\begin{theorem} \label{thm:extend-isometric-columns-to-unitary-matrix} \hlt{Extending Isometric Columns to a Unitary Matrix}\\
If $X\in\operatorname{Mat}(m,n,\mathbb{F})$ and $X^*X=I_n$, then there exists $X'\in\operatorname{Mat}(m,m-n,\mathbb{F})$ such that $[X\ X']$ is unitary.
\begin{proof}
Write $X=[\mathbf{x}_1\ \cdots\ \mathbf{x}_n]$. Since $(X^*X)_{ij}=\innerprod[\mathbb{F}^m]{\mathbf{x}_j,\mathbf{x}_i}$, the condition $X^*X=I_n$ says $\mathbf{x}_1,\ldots,\mathbf{x}_n$ are orthonormal in $\mathbb{F}^m$. By Corollary \ref{cor:orthonormal-list-extends-to-basis}, extend this list to an orthonormal basis $\mathbf{x}_1,\ldots,\mathbf{x}_m$ of $\mathbb{F}^m$.\\
Let $X'=[\mathbf{x}_{n+1}\ \cdots\ \mathbf{x}_m]$. Then the columns of $[X\ X']$ are orthonormal, so $[X\ X']$ is unitary.\\
\hlt{Structural Commentary}: Isometric columns are a partial orthonormal basis.\\
\hlt{Key Insight}: Complete the columns by extending the orthonormal list.
\end{proof}
\end{theorem}

\begin{definition} \label{def:unitary-and-permutation-similarity} \hlt{Unitary and Permutation Similarity}\\
Square matrices $A,B\in\operatorname{Mat}(n,n,\mathbb{F})$ are \hlt{unitarily similar} if $A=UBU^*$ for some unitary $U$. They are \hlt{orthogonally similar} if $\mathbb{F}=\mathbb{R}$ and $U$ can be chosen real orthogonal.\\
A \hlt{permutation matrix} is a square matrix with exactly one entry equal to $1$ in each row and each column, and all other entries equal to $0$. Matrices are \hlt{permutation similar} if $A=PBP^*$ for some permutation matrix $P$.
\end{definition}

\begin{proposition} \label{prop:permutation-matrices-are-orthogonal} \hlt{Permutation Matrices Are Orthogonal}\\
If $P=[\mathbf{e}_{\sigma(1)}\ \cdots\ \mathbf{e}_{\sigma(n)}]$ is a permutation matrix, then $P^*P=PP^*=I_n$. Hence $P$ is real orthogonal.
\begin{proof}
The columns of $P$ are a rearrangement of the standard orthonormal basis of $\mathbb{R}^n$, hence they are orthonormal. Theorem \ref{thm:unitary-matrix-equivalent-conditions} gives $P^*P=PP^*=I_n$.\\
\hlt{Structural Commentary}: A permutation matrix only reorders the standard basis.\\
\hlt{Key Insight}: Reordering orthonormal vectors preserves orthonormality.
\end{proof}
\end{proposition}

\begin{definition} \label{def:householder-matrix} \hlt{Householder Matrix}\\
Let $\mathbf{w}\in\mathbb{F}^n$ be nonzero and set $\mathbf{u}=\mathbf{w}/\norm{\mathbf{w}}_2$.\\
The \hlt{Householder matrix} associated to $\mathbf{w}$ is $U_{\mathbf{w}}\defeq I_n-2P_{\mathbf{u}}=I_n-2\frac{\mathbf{w}\mathbf{w}^*}{\norm{\mathbf{w}}_2^2}$.\\
It acts by $\mathbf{x}\mapsto U_{\mathbf{w}}\mathbf{x}=\mathbf{x}-2\innerprod[\mathbb{F}^n]{\mathbf{x},\mathbf{u}}\mathbf{u}$.
\end{definition}

\begin{figure}[H]
  \centering
  \begin{tikzpicture}[
  >=Stealth,
  every node/.style={font=\small},
  plane/.style={fill=gray!18, draw=gray!45, line width=0.7pt},
  normal/.style={->, black, line width=0.9pt},
  vec/.style={->, blue!75!black, line width=0.95pt},
  refl/.style={->, teal!75!black, line width=0.95pt},
  guide/.style={gray!60, dashed, line width=0.65pt}
]
  \coordinate (O) at (0,0);
  \coordinate (A) at (-2.2,-0.45);
  \coordinate (B) at (2.6,-0.45);
  \coordinate (C) at (1.75,0.55);
  \coordinate (D) at (-3.05,0.55);
  \coordinate (X) at (1.45,2.0);
  \coordinate (PX) at (0.72,0.05);
  \coordinate (RX) at (0.0,-1.9);

  \draw[plane] (A) -- (B) -- (C) -- (D) -- cycle;
  \node at (2.35,0.35) {$\mathbf{w}^{\perp}$};

  \draw[normal] (O) -- (0,2.2) node[above] {$\mathbf{w}$};
  \draw[normal] (O) -- (0,-1.05) node[left] {$-\mathbf{w}$};

  \draw[vec] (O) -- (X) node[above right] {$\mathbf{x}$};
  \draw[refl] (O) -- (RX) node[below] {$U_{\mathbf{w}}\mathbf{x}$};
  \draw[guide] (X) -- (PX);
  \draw[guide] (RX) -- (PX);
  \draw[gray!70] ($(PX)+(-0.14,0.08)$) -- ($(PX)+(0.02,0.2)$) -- ($(PX)+(0.16,0.08)$);
\end{tikzpicture}

  \caption{A Householder matrix fixes $\mathbf{w}^{\perp}$ and flips the $\operatorname{span}(\mathbf{w})$ component.}
  \label{fig:householder-reflection}
\end{figure}

\begin{theorem} \label{thm:householder-matrix-properties} \hlt{Householder Matrices Are Unitary Reflections}\\
Every Householder matrix $U_{\mathbf{w}}$ is unitary, Hermitian, and involutive: $U_{\mathbf{w}}^*U_{\mathbf{w}}=I_n$, $U_{\mathbf{w}}^*=U_{\mathbf{w}}$, $U_{\mathbf{w}}^2=I_n$.\\
If $\mathbf{w}\in\mathbb{R}^n$, then $U_{\mathbf{w}}$ is real orthogonal and symmetric.
\begin{proof}
Let $\mathbf{u}=\mathbf{w}/\norm{\mathbf{w}}_2$. By Proposition \ref{prop:rank-one-projection-matrix-properties}, $P_{\mathbf{u}}^*=P_{\mathbf{u}}$ and $P_{\mathbf{u}}^2=P_{\mathbf{u}}$.\\
Therefore $U_{\mathbf{w}}^*=I_n-2P_{\mathbf{u}}=U_{\mathbf{w}}$, $U_{\mathbf{w}}^2=(I_n-2P_{\mathbf{u}})^2=I_n-4P_{\mathbf{u}}+4P_{\mathbf{u}}^2=I_n$.\\
Thus $U_{\mathbf{w}}^*U_{\mathbf{w}}=U_{\mathbf{w}}^2=I_n$, so $U_{\mathbf{w}}$ is unitary. The real case is the same statement with conjugation removed.\\
\hlt{Structural Commentary}: A Householder matrix reflects across an orthogonal hyperplane.\\
\hlt{Key Insight}: The identity $P_{\mathbf{u}}^2=P_{\mathbf{u}}$ makes $(I-2P_{\mathbf{u}})^2=I$.
\end{proof}
\end{theorem}

\begin{theorem} \label{thm:householder-maps-equal-norm-vectors} \hlt{Householder Map Between Equal-Norm Vectors}\\
Let $\mathbf{x},\mathbf{y}\in\mathbb{F}^n$ satisfy $\norm{\mathbf{x}}_2=\norm{\mathbf{y}}_2\neq0$. Define $\sigma=
\begin{cases}
1, & \innerprod[\mathbb{F}^n]{\mathbf{x},\mathbf{y}}=0,\\
-\dfrac{\overline{\innerprod[\mathbb{F}^n]{\mathbf{x},\mathbf{y}}}}{\abs{\innerprod[\mathbb{F}^n]{\mathbf{x},\mathbf{y}}}}, & \innerprod[\mathbb{F}^n]{\mathbf{x},\mathbf{y}}\neq0,
\end{cases}$, where $\mathbf{w}=\mathbf{y}-\sigma\mathbf{x}$.\\
Then $\mathbf{w}\neq\mathbf{0}$, $\abs{\sigma}=1$, and $\sigma U_{\mathbf{w}}\mathbf{x}=\mathbf{y}$.
\begin{proof}
Clearly $\abs{\sigma}=1$. Let $z=\innerprod{\mathbf{x},\mathbf{y}}$ and $r=\norm{\mathbf{x}}_2=\norm{\mathbf{y}}_2$. By construction, $\sigma z=-\abs{z}$, with the convention that this equals $0$ when $z=0$. Hence $\norm{\mathbf{w}}_2^2 =\norm{\mathbf{y}-\sigma\mathbf{x}}_2^2 =2r^2-\sigma z-\overline{\sigma z} =2\left(r^2+\abs{z}\right)>0$, so $\mathbf{w}\neq\mathbf{0}$.\\
Since $\mathbf{y}=\mathbf{w}+\sigma\mathbf{x}$, $\sigma U_{\mathbf{w}}\mathbf{x} =\sigma\mathbf{x}-2\sigma\frac{\innerprod{\mathbf{x},\mathbf{w}}}{\norm{\mathbf{w}}_2^2}\mathbf{w} =\mathbf{y}-\left(1+2\sigma\frac{\innerprod{\mathbf{x},\mathbf{w}}}{\norm{\mathbf{w}}_2^2}\right)\mathbf{w}$.\\
But $\innerprod{\mathbf{x},\mathbf{w}}=z-\overline{\sigma}r^2$, so $2\sigma\innerprod{\mathbf{x},\mathbf{w}}=2\sigma z-2r^2=-2(r^2+\abs{z})=-\norm{\mathbf{w}}_2^2$.\\
Coefficient of $\mathbf{w}$ is hence zero, proving $\sigma U_{\mathbf{w}}\mathbf{x}=\mathbf{y}$.\\
\hlt{Structural Commentary}: The scalar $\sigma$ chooses the phase so that $\mathbf{y}-\sigma\mathbf{x}$ gives a stable reflection direction.\\
\hlt{Key Insight}: Householder reflections connect equal-norm vectors up to a unit scalar.
\end{proof}
\end{theorem}

\begin{corollary} \label{cor:householder-unitary-annihilator} \hlt{Householder Unitary Annihilator}\\
Let $\mathbf{x}=[x_1\ \cdots\ x_n]^t\in\mathbb{F}^n$ be nonzero. Define $\sigma=
\begin{cases}
1, & x_1=0,\\
-\dfrac{\overline{x_1}}{\abs{x_1}}, & x_1\neq0,
\end{cases}$, where $\mathbf{w}=\norm{\mathbf{x}}_2\mathbf{e}_1-\sigma\mathbf{x}$.\\
Then $\sigma U_{\mathbf{w}}\mathbf{x}=\norm{\mathbf{x}}_2\mathbf{e}_1$. If $\mathbf{x}$ is a unit vector, then $\overline{\sigma}U_{\mathbf{w}}$ is unitary and has $\mathbf{x}$ as its first column.
\begin{proof}
Apply Theorem \ref{thm:householder-maps-equal-norm-vectors} with $\mathbf{y}=\norm{\mathbf{x}}_2\mathbf{e}_1$, noting that $\innerprod{\mathbf{x},\mathbf{y}}=x_1\norm{\mathbf{x}}_2$. If $\norm{\mathbf{x}}_2=1$, then $(\sigma U_{\mathbf{w}})\mathbf{x}=\mathbf{e}_1$, so $(\sigma U_{\mathbf{w}})^*\mathbf{e}_1=\mathbf{x}$. Since $(\sigma U_{\mathbf{w}})^*=\overline{\sigma}U_{\mathbf{w}}$, this says the first column of $\overline{\sigma}U_{\mathbf{w}}$ is $\mathbf{x}$.\\
\hlt{Structural Commentary}: The reflector concentrates all mass of $\mathbf{x}$ into the first coordinate.\\
\hlt{Key Insight}: Householder matrices produce explicit unitary matrices with a prescribed first column.
\end{proof}
\end{corollary}
